\documentclass[10pt]{article}

\usepackage[utf8]{inputenc}

\usepackage[T1]{fontenc}

\usepackage{amsmath}

\usepackage{amsfonts}

\usepackage{amssymb}

\usepackage[version=4]{mhchem}

\usepackage{stmaryrd}

\usepackage{bbold}

\usepackage{mathrsfs}



\begin{document}

свойством о т с у т с т в и я п о с е д е й с т в и я. Оно является математич. выражением интуитивного представления об отсутствии «памяти» о прошлом в развитии процесса. Таким образом, М. п.- это процесс, обладающий марковским свойством.



Имеются отличающиеся друг от друга точные математич. определения марковского свойства и тем самым М. П. Наиболее употребительным в математич. физике является следующее. Пусть $(X, \mathfrak{B})$ - измеримое фазовое пространство процесса и пусть $\mathfrak{F}_{t}$ (либо $\mathfrak{F}^{t}$ ) $-\sigma$-алгебра событий, порожденная всеми событиями вида $\{\xi(s) \in \Gamma\}, s \in T, \Gamma \in \mathfrak{B}, s \leqslant t$ (либо $s \geqslant t$ ). Тогда марковское свойство состоит в том, что для всех $t, A_{1} \in \mathfrak{F}_{t}, A_{2} \in \mathfrak{F}^{t}$ и почти наверное для всех значений $\xi(t)$ выполняется равенство



\[

\mathrm{P}\left\{A_{1} A_{2} \mid \xi(t)\right\}=\mathrm{P}\left\{A_{1} \mid \xi(t)\right\} \mathrm{P}\left\{A_{2} \mid \xi(t)\right\}

\]



или, что то же самое, для любого $t$ и $A \in \mathfrak{F}^{t}$ с вероятностью 1 справедлива формула



\[

\mathrm{P}\left\{A \mid \widetilde{\mho}_{t}\right\}=\mathrm{P}\{A \mid \xi(t)\} .

\]



Это определение является непосредственным обобщением марковского свойства, отвечающего Маркова цепи - М. П. с дискретным временем. Чтобы отличить его от других не эквивалентных ему определений, иногда процесс $\{\xi(t)\}$, обладающий свойством (1) либо (2), называют м а р к о в ской случай ной функцие й (точнее не обры ва ю щейся марковской случайной функцие й).



Будем полагать, что М. П. задан на $\left[t_{0}, \infty\right)$. Следствием марковского свойства (2) и формулы полной вероятности является соотношение



$\mathrm{P}\left\{\xi(t) \in \Gamma \mid \xi\left(t^{\prime}\right)\right\}=\mathrm{E}\left\{\mathrm{P}\{\xi(t) \in \Gamma \mid \xi(s)\} \mid \xi\left(t^{\prime}\right)\right\}, t \geqslant s \geqslant t^{\prime}$,



справедливое при всех $\Gamma \in \mathfrak{B}$ и почти наверное для $\xi\left(t^{\prime}\right)$ относительно меры, связанной с процессом. Это соотношение называется уравнением Колмогорова-Чепм е на. При моделировании случайной эволюции системы в задачах математич. физики на основе марковского процесса фазовое пространство системы $(X, \mathfrak{B})$, как правило, изоморфно борелевскому подмножеству полного сепарабельного метрич. пространства с $\sigma$-алгеброй его борелевских подмножеств (в большинстве случаев $X \subset \mathbb{R}^{m}$ ). В этом случае для М. п. существует условная пе реходная в е роятность, или стохастическое ядро, то есть такая функция $P(s, x ; t, \Gamma), x \in X, \Gamma \in \mathfrak{B}, s<t$, к-рая:



\begin{enumerate}

  \item при фиксированных $s, x, t$ является вероятностной мерой на $(X, \mathfrak{B})$;



  \item $\mathfrak{B}$ - измерима при фиксированных $s, t$, $\Gamma$;



  \item с вероятностью 1 при всех $s, t$, Г



\end{enumerate}



\[

P(s, \xi(s) ; t, \Gamma)=\mathrm{P}\{\xi(t) \in \Gamma \mid \xi(s)\} .

\]



Кроме того, в указанной ситуации уравнение Колмогорова-Чепмена, справедливое при всех $x \in X, \Gamma \in \mathfrak{B}$ и $t>t^{\prime}$, для переходных вероятностей записывается в виде



\[

P\left(t^{\prime}, x ; t, \Gamma\right)=\int_{X} P(s, y ; t, \Gamma) P\left(t^{\prime}, x ; s, d y\right) .

\]



Стохастич. ядро вместе с мерой



\[

\mu(\Delta)=\mathrm{P}\left\{\xi\left(t_{0}\right) \in \Delta\right\}, \Delta \in \mathfrak{B} \text {, }

\]



определяют все конечномерные распределен и я М. П. и, следовательно, определяют М.П. с точностью до с тохастичес кой эк в и вале н т ности, то есть М. П. в широком смысле:



$P\left\{\xi\left(t_{i}\right) \in \Delta_{i} ; i=1, \ldots, n\right\}=$



\[

\begin{gathered}

=\int_{X} \mu\left(d x_{0}\right) \int_{\Delta_{1}} P\left(t_{0}, x_{0} ; t_{1}, d x_{1}\right) \ldots \int_{\Delta_{n}} P\left(t_{n-1}, x_{n-1} ; t_{n}, d x_{n}\right), \\

\Delta_{1}, \Delta_{2}, \ldots, \Delta_{n} \in \mathfrak{B} ; t_{0}<t_{1}<\ldots<t_{n} .

\end{gathered}

\]



Иногда приходится рассматривать системы, для к-рых определение М. П. с помощью стохастич. ядра недостаточно. Такая ситуация имеет место, напр., для рождения и гибели процессов, когда с ненулевой вероятностью выбороч- ные функции процесса покидают фазовое пространство системы - происходит исчезновение процесса. Такие случайные процессы называются о 6 р ы в а ю и м и с я, и с их помощью моделируется случайная эволюция, приводящая к таким физич. явлениям, как обрыв, взрыв, гибель и т. Д. Вероятностное пространство $\{\Omega, \mathscr{F}, P\}$ обрывающегося процесса строится по следующей схеме. Фазовое пространство процесса расширяется за счет добавления к $X$ элемента $\mathfrak{i}$ и расширения $\sigma$-алгебры $\mathfrak{B}$ до $\sigma$-алгебры $\mathfrak{B}_{\mathfrak{i}}$; путем присоединения к ней множеств вида $A \cup\{\mathfrak{i}\}, A \in \mathfrak{B}$. Тогда пространство $\Omega$ реализаций процесса состоит из функций $\xi(t)$ со значениями в $X \cup\{i\}$, а $\sigma$-алгебра $\mathfrak{F}_{t}\left(\lim \mathscr{F}_{t}=\mathscr{F}\right)$ порождается всеми событиями $\{\xi(s) \in \Gamma\}, \Gamma \in \mathfrak{B}_{\mathfrak{i}}, s \leqslant t$. Случайное время обрыва процесса $\tau(\xi)$ в этом случае равно



\[

\tau(\xi)=\inf \{t: \xi(t)=\mathfrak{i}\}

\]



Обрывающийся процесс $\{\xi(t)\}$, удовлетворяющий марковскому свойству (2), называется о б ры в а ю и м с я М. П. (о 6 ры в а ю е йс я марковской случайной функцией). При тех же предположениях относительно фазового пространства $(X, \mathfrak{B})$, что и в случае необрывающегося М. П., существует условная переходная вероятность $P(s, x ; t, \Gamma)$, $x \in X \cup\{i\}$, к-рая обладает свойствами 1)-3) и удовлетворяет уравнению (3). При $x=i$ переходная вероятность равна 1 , если $\mathfrak{i} \in \Gamma$, и 0 , если $\mathfrak{i} \notin \Gamma$. При $x \neq \mathfrak{i}$ и Г $\subset X$ функция $P(s, x ; t, \Gamma)$ называется полустохастическим ядром М. П. $\{\xi(t)\}$, так как $P(s, x ; t, X)<1$. При этом $[1-P(s, x ; t, X)]$ есть вероятность обрыва процесса к моменту $t$ при условии, что $\check{\zeta}(s)=x$.



Примерами М.П. могут служить винеровский процесс, пуассоновский процесс и т. Д., К-рые находят широкое применение при моделировании случайной эволюции физич. систем.



М. П. являются также решения стохастических дифференциальных уравнений по винеровскому процессу.



Лит.: [1] Гихман И. И., Ско роход А. В., Теория случайных процессов, т. 2, М., 1973; [2] Д ы н к и н Е. Б., Марковские процессы, М., 1963; [3] В е н т це л ь А. Д., Курс теории случайных процессов, М., 1975. Ю. П. Вирченко. МАРКОВСКИЙ СЛУЧАЙНЫЙ ПРОЦЕСС - см. Марковское случайное поле.



MAPKOBCKOE OTOБPAXEHИE, д и н а м и ч е с к о е о т ображе ни е, - отображение, описывающее конечное изменение наблюдаемых или состояний системы в результате, вообще говоря, необратимой эволюции в алгебраической формулировке статистической механики, а также в некоммутативной теории вероятностей. Пусть $\mathfrak{A}$ и $\mathfrak{B}-$ нек-рые $C^{*}$-алгебры с единицами. П е р е хо д ны м о то6 ражением называется вполне положительное отображение $\Phi$ из $\mathfrak{A}$ в $\mathfrak{B}$, переводящее единицу из алгебры $\mathfrak{U}$ в единицу из алгебры $\mathfrak{B}$. Если $\mathfrak{A}=\mathfrak{B}$, то переходное отображение называется ма к ко в ским от о бражением.



Для всякого переходного отображения существует сопряженное $\Phi^{*}$, к-рое аффинно и слабо *-непрерывно отображает множество состояний $\mathscr{S}_{\mathfrak{B}}$ в $\mathfrak{S}_{\mathfrak{A}}$. Если $\mathfrak{A}$ и $\mathfrak{B}$ являются $W^{*}$-алгебрами, а отображение нормально в том смысле, что



\[

\sup _{\alpha} \Phi\left(X_{\alpha}\right)=\Phi\left(\sup _{\alpha} X_{\alpha}\right)

\]



для любого ограниченного направленного множества самосопряженных элементов $\left\{X_{\alpha}\right\} \subset \mathfrak{A}$, то $\Phi^{*}$ переводит нормальные состояния в нормальные же.



П р и м е р ы. 1) Пусть $\mathfrak{A}=\mathfrak{B}-$ коммутативная алгебра ограниченных измеримых функций на измеримом пространстве $\Omega$ и $\mathrm{P}\left(d \omega \mid \omega^{\prime}\right)$ - переходная вероятность в $\Omega$. Соотношение



\[

\Phi(X)\left(\omega^{\prime}\right)=\int_{\Omega} \mathrm{P}\left(d \omega \mid \omega^{\prime}\right) X(\omega), \quad X \in \mathfrak{A},

\]



определяет М. о. При этом $\Phi^{*}$ порождает аффинное отображение в выпуклом множестве всех вероятностных мер на $\Omega$.



\begin{enumerate}

  \setcounter{enumi}{1}

  \item Пусть $\mathfrak{A}=\mathfrak{B}=\mathfrak{B}(H)$ - алгебра всех ограниченных операторов в гильбертовом пространстве $H$. Всякое нормальное

\end{enumerate}



\end{document}